\section{Experiments}
\subsection{Experimental Setup}
Following Flow-GRPO~\cite{liu2025flow}, we evaluate our method on four tasks: GenEval~\cite{geneval}, OCR~\cite{textdiffuser}, PickScore~\cite{pickscore}, and DeQA~\cite{deqa}. We adopt the official checkpoints as expert teachers for the first three tasks. The DeQA teacher is specifically trained across the three datasets by blending DeQA and PickScore rewards at a 4:6 ratio. All training and test data strictly follow the Flow-GRPO splits. Training is executed on 4 nodes ($8 \times \text{H800}$ GPUs each), while evaluation is conducted on a single $8 \times \text{H800}$ node.

We primarily evaluate Flow-OPD against two categories of baselines: 
(1) \textbf{Monolithic-Reward GRPO}, denoted as \textit{GRPO-[reward name]}, where the model is fine-tuned using Flow-GRPO on a single reward objective; 
(2) \textbf{Hybrid-Reward GRPO}, denoted as \textit{GRPO-Mix}, which employs a weighted reward combination with a fixed ratio of \textit{GenEval} : \textit{OCR} : \textit{PickScore} = 3 : 1 : 1. 
These baselines serve to highlight the limitations of conventional scalar-based alignment when scaling to multi-dimensional expert capabilities.
%我们主要评估flow-opd在四类任务上的表现:1.图像合成-genevak，2.视觉文本渲染，3.图像生成质量（deqa）和4.人类偏好一致性（pickscore）。
%教师训练：对于图像合成，文本渲染和pickscore一致性，我们直接使用flow-grpo提供的官方检查点，对于deqa，我们使用flow-grpo的官方仓库进行复现。详细的prompt和reward计算方法请参考补充材料。

% To evaluate the performance of \textbf{Flow-OPD}, we conduct experiments across four key categories of tasks. To ensure a fair and consistent comparison, all training prompts and their respective distributions strictly follow the protocols established by \textbf{Flow-GRPO}. Detailed hyperparameters and compute specifications are provided in .

% \paragraph{Compositional Image Generation (GenEval)} This task assesses the model's ability to handle complex semantic compositions, such as object counting, spatial relations, and attribute binding. We utilize the \textbf{GenEval} benchmark and adopt the specific prompt weighting from Flow-GRPO, where the ratio is set as $Position : Counting : Attribute Binding : Colors : Two Objects : Single Object = 7 : 5 : 3 : 1 : 1 : 0$. This distribution prioritizes the most challenging aspects of spatial and numerical reasoning.

% \paragraph{Visual Text Rendering} Accurate text synthesis within images is a critical metric for T2I models. Following the setup in Flow-GRPO, we utilize 20K training prompts generated by GPT-4o. Each prompt follows a fixed template: \textit{``A sign that says `text' ''}, requiring the model to render the exact string within the visual context. This task evaluates the model's precision and character-level alignment.

% \paragraph{Image Quality Evaluation (DeQA)} To mitigate the risk of \textit{reward hacking}—where the model might optimize for task-specific rewards at the expense of overall visual fidelity—we monitor image quality using \textbf{DeQA}. Following Flow-GRPO, we perform evaluations on the \textbf{DrawBench} dataset to ensure that our method maintains high aesthetic standards and structural integrity across diverse and open-domain prompts.

% \paragraph{Human Preference Alignment (PickScore)} This task aims to align the generated outputs with general human aesthetic and intent preferences. We utilize \textbf{PickScore}, a metric derived from large-scale human-annotated pairwise comparisons, to evaluate the model's capacity to generate images that are both semantically accurate and visually superior according to human-centric benchmarks.
\subsection{Main Results}

\begin{figure*}[t]
    \centering
    \includegraphics[scale=0.49]{images/p1.pdf}
    % \vspace{-15pt}
\caption{Qualitative comparison between Flow-OPD and various baselines across diverse tasks. Our method consistently demonstrates superior instruction-following capabilities, delivering high-fidelity image synthesis and structural coherence that align more closely with human preferences.}
    \label{fig:main}
    % \vspace{-0.4cm}
\end{figure*}

\input{tables/main_result_0428}

The quantitative results in Table~\ref{tab:model_comparison_updated} demonstrate that Flow-OPD consistently matches or surpasses the specialized teacher models across all benchmarks, particularly in text rendering and DeQA image quality. Crucially, it resolves the severe cross-domain interference inherent to specialization (e.g., the PickScore teacher's GenEval score dropping to 0.51) and overcomes the optimization bottlenecks of sparse-reward multi-task GRPO. By leveraging dense multi-expert supervision, Flow-OPD seamlessly consolidates diverse expertise without capability degradation.

Qualitative results in Fig.~\ref{fig:main}) show that Flow-OPD achieves an optimal multi-task trade-off, balancing high prompt fidelity with superior visual aesthetics. Remarkably, Flow-OPD succeeds in certain edge cases where all individual teachers fail, a phenomenon we term \textit{Teacher-Surpassing}. We hypothesize this emergent superiority stems from knowledge cross-pollination within the latent flow manifold. While individual teachers are constrained by domain-specific biases, simultaneous dense guidance forces the student to learn a more holistic, smoothed representation. This collective supervision bridges epistemic gaps, enabling the student to synthesize novel trajectories that ultimately surpass any single supervisor.
% The quantitative results, as illustrated in Table ~\ref{tab:model_comparison}, demonstrate that Flow-OPD achieves highly competitive performance across all four benchmarks, consistently reaching the high-water marks set by the respective teacher models. Notably, our framework slightly surpasses the specialized teachers in both text rendering and deqa image quality, underscoring the efficacy of our proposed methodology in consolidating and even enhancing domain-specific expertise.

% A closer inspection underscores the cross-domain interference present during specialization. Specifically, the PickScore teacher exhibits a marked performance degradation in GenEval (dropping to 0.51), which underscores the necessity of maintaining generalist capabilities. Additionally, multi-task reinforcement learning via GRPO encounters a optimization bottleneck due to discrete reward landscapes, failing to match the proficiency of task-specific teachers. By leveraging dense supervisory signals from multiple experts and a disentangled optimization objective, Flow-OPD mitigates these conflicts, manifesting powerful capabilities across all evaluation metrics that rival the state-of-the-art specialized teachers.

% Qualitative results, as depicted in Fig.~\ref{fig:main}, illustrate that Flow-OPD achieves an optimal multi-task trade-off. On one hand, our model demonstrates high prompt fidelity, successfully executing complex instructions involving specific attributes, object relations, and even the challenging task of case-sensitive text rendering. Furthermore, it exhibits superior visual aesthetics that align closely with human preference. Remarkably, we observe that Flow-OPD excels in certain edge cases where all individual teacher models fail—a phenomenon we term 'Teacher-Surpassing'. 
% We hypothesize that this emergent superiority stems from knowledge cross-pollination within the latent flow manifold. While individual teachers are often constrained by domain-specific biases or local minima in their respective optimization landscapes, the student model, under the simultaneous guidance of multiple experts, is forced to learn a more holistic and smoothed representation. This collective supervision effectively 'fills' the epistemic gaps of single teachers, enabling the student to interpolate and synthesize novel trajectories that reside in the intersection of expert competencies, ultimately surpassing the capabilities of any single supervisor.

%几个实验
%1. sft的必要性
%2. 
\subsection{Analysis}
\subsubsection{Cold Start Ablation}

\begin{figure*}[t]
    \centering
    \includegraphics[scale=0.58]{images/bar.pdf}
    % \vspace{-15pt}
    \caption{Cold-start ablation results.
    }
    \label{fig:cold—start}
    % \vspace{-0.4cm}
\end{figure*}

\begin{figure*}[t]
    \centering
    \includegraphics[scale=0.445]{images/p2.pdf}
    % \vspace{-15pt}
    \caption{Qualitative ablation results of Manifold Anchor Regularization.
    }
    \label{fig: kl}
    % \vspace{-0.4cm}
\end{figure*}

\input{tables/compbench}

\input{tables/kl}

As shown in Fig.~\ref{fig:cold—start}, cold-start initialization rapidly establishes a robust foundation for subsequent training. Between the two regimes, Supervised Fine-Tuning (SFT) serves as a widely adopted and highly scalable strategy; notably, its inherent flexibility presents a promising avenue for extracting capabilities from heterogeneous teachers in future applications. Conversely, model merging optimally leverages the available homogeneous teachers for superior functional alignment without any additional training costs. 
%
Crucially, Flow-OPD consistently outperforms both from-scratch and cold-started multi-task GRPO. While GRPO converges to sub-optimal states due to inter-task conflicts caused by sparse scalar rewards, Flow-OPD leverages dense multi-expert supervision to resolve gradient interference. Consequently, our method achieves substantial, uniform gains across all baselines, successfully matching or exceeding the performance ceilings of individual specialized teachers.

\subsubsection{OOM Generalization}

To further investigate the generalization capabilities of our method, we conduct additional evaluations on the T2I-CompBench benchmark. Flow-OPD demonstrates superior out-of-domain generalization compared to multi-task GRPO, achieving state-of-the-art (SOTA) performance across multiple compositional metrics. Notably, when initialized from the identical cold-start baseline, standard GRPO suffers from catastrophic forgetting in specific capability dimensions, such as shape rendering and 3D spatial relations. In contrast, by leveraging dense multi-expert supervision and task-style decoupling regularization, Flow-OPD effectively mitigates these regression issues, yielding robust and comprehensive performance enhancements.

\subsubsection{Manifold Anchor Regularization}


Manifold Anchor Regularization (MAR) is a task-agnostic constraint designed to maintain generative integrity and aesthetic alignment. As shown in Fig.~\ref{fig: kl}, vanilla GRPO-based optimization often triggers background mode collapse—where models overfit to monotonous environments—and semantic redundancy, leading to identical features across multiple entities due to coarse reward granularity. While teachers like DeQA provide diverse samples, they often struggle with instruction following. MAR resolves these issues by anchoring optimization to a high-fidelity manifold, balancing structural diversity with precise semantic adherence.
Table~\ref{tab:reward_metrics} further provides quantitative evidence of our method's superiority in image quality and human preference alignment. The integration of MAR leverages additional supervision across the entire dataset, significantly enhancing both the visual quality and the expressive power of the generated images.